\section{Introduzione}
\label{sec:introduction}

Il Federated Learning permette a un insieme di nodi client, ciascuno in
possesso di un dataset locale privato, di addestrare collaborativamente un
modello condiviso senza scambiare i dati grezzi. Nella formulazione centralizzata,
un server raccoglie i pesi locali
di ogni client, li aggrega secondo un particolare algoritmo e ridistribuisce il modello aggiornato: il server è
un single point of failure, la sua banda è un collo di bottiglia che cresce
con il numero di client, e ogni round richiede una sincronizzazione globale
al passo del client più lento.

Il presente lavoro affronta questi limiti progettando un sistema di FL
\emph{decentralizzato}: i pesi vengono scambiati tra i client tramite un
protocollo di comunicazione peer-to-peer a push one-hop e fanout fisso, in
cui ciascun client invia periodicamente i propri pesi a un sottoinsieme di
vicini scelti casualmente, e vengono aggregati localmente da chi li riceve
tramite FedAvg. Un componente centralizzato, il discovery server,
resta limitato esclusivamente alla registrazione e alla scoperta reciproca dei
client.
